\chapter{Privacy-free Free-XOR garbling}
\label{chap:garbling}

Labels live in a vector space $V$ over $\F$; for instance $V = \F^{128}$, that is,
128-bit strings with XOR. A wire carries one of two labels, $L$ for the bit $0$ and $L + \Delta$
for the bit $1$, where $\Delta \in V$ is a global offset chosen by the garbler. The evaluator
holds one of them, the \emph{active} label; the other one is \emph{inactive}.

\section{Circuits and plaintext evaluation}

\begin{definition}[Circuits]
  \label{def:circuit}
  \lean{GcZk.Node, GcZk.Circuit}
  \leanok
  A node is one of $\mathsf{input}(j)$, $\mathsf{const}(b)$, $\mathsf{xor}(a, b)$ and
  $\mathsf{mux}(s, z, o)$, where $j, a, b, s, z, o \in \N$ and $b \in \{0, 1\}$. A circuit is a
  list of nodes. The operands are positions of earlier nodes; a position that does not refer to
  an earlier node reads the default wire (value $0$, label $0$).
\end{definition}

The multiplexer $\mathsf{mux}(s, z, o)$ outputs the value of $o$ if $s$ is $1$ and the value
of $z$ otherwise; AND and OR are multiplexers with a constant candidate. Accepting arbitrary
positions avoids carrying a well-formedness hypothesis through every statement.

\begin{definition}[Plaintext evaluation]
  \label{def:plain}
  \lean{GcZk.plainNode, GcZk.plain}
  \leanok
  \uses{def:circuit}
  For an input $x \colon \N \to \{0, 1\}$, the values of the nodes of $C$ are computed in order:
  $x_j$, $b$, $v_a \oplus v_b$, and $v_o$ or $v_z$ according to $v_s$.
\end{definition}

\section{Garbling and evaluation}

\begin{definition}[Labels of a bit]
  \label{def:lab}
  \lean{GcZk.lab}
  \leanok
  For a zero-label $L$ and a bit $b$, $\lab_\Delta(L, b) = L + b\,\Delta$.
\end{definition}

\begin{definition}[Garbling]
  \label{def:garble}
  \lean{GcZk.Seed, GcZk.Material, GcZk.garbleNode, GcZk.garble}
  \leanok
  \uses{def:circuit, def:lab}
  The garbler's seed consists of $\Delta$ and the zero-labels of the inputs and of the
  constants; the hash is a map $H \colon \N \times V \to V$ (gate position, label). The node at
  position $i$ receives a zero-label $L_i$ and publishes some material:
  \begin{itemize}
    \item $\mathsf{input}(j)$: $L_i$ is the zero-label of input $j$; no material.
    \item $\mathsf{const}(b)$: $L_i$ is the zero-label of the constant; the material is its
    active label $\lab_\Delta(L_i, b)$.
    \item $\mathsf{xor}(a, b)$: $L_i = L_a + L_b$ (Free XOR); no material.
    \item $\mathsf{mux}(s, z, o)$: $L_i = H(i, L_s) + L_z$, and the material is the single
    ciphertext
    \[ T_i = H(i, L_s) + H(i, L_s + \Delta) + L_z + L_o . \]
  \end{itemize}
\end{definition}

\begin{remark}[The multiplexer is a half-gate]
  \label{rem:mux-half-gate}
  Since $\mathsf{mux}(s, z, o) = z \oplus (s \wedge (z \oplus o))$, the multiplexer is an AND
  gate whose first input $s$ is known to the evaluator, followed by a free XOR with $z$. The
  evaluator half-gate of Zahur, Rosulek and Evans \cite[Sect.~3.1]{ZRE15} garbles $c = a \wedge b$,
  with $a$ known to the evaluator and zero-labels $A$, $B$, by choosing the zero-label $C = H(A)$
  and sending the single ciphertext $H(A) \oplus H(A \oplus R) \oplus B$, where $R$ is the global
  offset. With $A = L_s$, $B = L_z + L_o$ (the zero-label of $z \oplus o$) and $R = \Delta$, this is
  $T_i$, and the free XOR with $z$ turns $C$ into $L_i = H(i, L_s) + L_z$. The gate position $i$ in
  the hash plays the role of their gate index $j$ \cite[Sect.~3.2]{ZRE15}. Garbling every AND gate this way, with free XOR
  gates, is their privacy-free scheme \cite[Sect.~6]{ZRE15}.
\end{remark}

\begin{definition}[Evaluation]
  \label{def:eval}
  \lean{GcZk.Wire, GcZk.evalNode, GcZk.eval, GcZk.activeInputs}
  \leanok
  \uses{def:garble}
  The evaluator knows the value and the active label of every wire. At a multiplexer with
  selector value $v_s$ and active selector label $S$, it takes the candidate selected by $v_s$,
  with active label $Y$, and outputs the label $H(i, S) + Y$, plus $T_i$ if $v_s = 1$. The active
  inputs for $x$ are the labels $\lab_\Delta(L_j, x_j)$.
\end{definition}

\section{Correctness}

\begin{lemma}[One step]
  \label{lem:good-snoc}
  \lean{GcZk.Good, GcZk.good_snoc}
  \leanok
  \uses{def:eval, def:plain}
  If every evaluated wire so far carries its plaintext value $v$ and the label
  $\lab_\Delta(L, v)$ of its zero-label $L$, the same holds after evaluating one more node.
\end{lemma}

\begin{proof}
  \leanok
  Only the multiplexer needs a computation. If $v_s = 0$ the evaluator holds $S = L_s$ and gets
  $H(i, L_s) + L_z + v_z \Delta = \lab_\Delta(L_i, v_z)$. If $v_s = 1$ it holds
  $S = L_s + \Delta$ and gets
  \[ H(i, L_s + \Delta) + L_o + v_o\Delta + T_i = H(i, L_s) + L_z + v_o \Delta , \]
  since every other term appears twice and $V$ has characteristic $2$.
\end{proof}

\begin{theorem}[Correctness]
  \label{thm:eval-garble}
  \lean{GcZk.eval_garble}
  \leanok
  \uses{lem:good-snoc}
  For every circuit $C$, seed, hash $H$ and input $x$, evaluating the garbling of $C$ on the active
  inputs of $x$ gives, at every node $i$, the plaintext value $v_i$ and the label
  $\lab_\Delta(L_i, v_i)$.
\end{theorem}

\begin{proof}
  \leanok
  Induction on the nodes, with Lemma~\ref{lem:good-snoc}.
\end{proof}

\begin{remark}
  No hypothesis on $H$ or on $\Delta$ is needed: correctness is pure algebra in characteristic
  $2$. It is the correctness of the scheme for an honest garbler.
\end{remark}
